\chapter{Proposta}
\label{sec:proposta}

\section{Requisitos}
    \subsection{Requisitos Funcionais}
        \begin{itemize}
            \item RF1: A aplicação deve permitir a execução de ferramentas de testes de penetração em um ambiente virtualizado.
            \item RF2: A aplicação deve oferecer uma interface gráfica intuitiva para facilitar o uso por profissionais iniciantes.
            \item RF3: A aplicação deve normalizar as saídas das ferramentas em um formato único para facilitar a análise dos resultados.
            \item RF4: A aplicação deve gerar relatórios técnicos detalhados com base nos resultados obtidos.
            \item RF5: A aplicação deve oferecer fluxos guiados para auxiliar os usuários durante a utilização das ferramentas.
        \end{itemize}

    \subsection{Requisitos Não Funcionais}
        \begin{itemize}
            \item RNF1: A aplicação deve ser desenvolvida em Python, utilizando a biblioteca Tkinter para a interface gráfica.
            \item RNF2: A aplicação deve utilizar Docker para virtualizar as ferramentas de testes de penetração, garantindo portabilidade e consistência entre diferentes ambientes de execução.
            \item RNF3: A aplicação deve ser compatível com os principais sistemas operacionais, incluindo Windows, Linux e macOS.
            \item RNF4: A aplicação deve ser projetada para ser fácil de usar, mesmo para usuários com pouca experiência em testes de penetração.
            \item RNF5: A aplicação deve ser segura, evitando a execução de comandos maliciosos e garantindo a integridade do sistema hospedeiro.
        \end{itemize}

\section{Arquitetura do Sistema}

\section{Ferramentas de Testes de Penetração}
    Ferramentas de testes de penetração são empregadas para apoiar a identificação, exploração e análise de vulnerabilidades
    em sistemas computacionais. Elas automatizam tarefas repetitivas, auxiliam na coleta de informações e permitem avaliar,
    de forma prática, a postura de segurança de redes, sistemas e aplicações. O uso dessas ferramentas
    é essencial para tornar o processo de pentest mais eficiente e sistemático, especialmente em ambientes complexos\citep{moreno2019introduccao}.

    A seguir, apresentam-se as principais ferramentas consideradas neste trabalho, classificadas de acordo com seu papel no processo de testes de penetração.

        \subsection{Nmap}

            O Nmap (Network Mapper) é uma ferramenta amplamente utilizada para reconhecimento e enumeração de redes. Seu principal
            objetivo é identificar hosts ativos, portas abertas, serviços em execução e possíveis sistemas operacionais. A fase de reconhecimento 
            é fundamental para o sucesso do pentest, pois fornece informações iniciais que orientam as etapas
            subsequentes de análise e exploração\citep{stallings2017network}.

        \subsection{Nuclei}

            O Nuclei é um framework de varredura baseado em templates, empregado para a detecção automatizada de vulnerabilidades conhecidas e más configurações.
            Seu modelo baseado em assinaturas facilita a identificação rápida de falhas recorrentes em aplicações e serviços, contribuindo para a eficiência da
            fase de enumeração de vulnerabilidades \citep{moreno2019introduccao}.
            
        \subsection{Nikto}

            O Nikto é um scanner de servidores web utilizado para identificar configurações inseguras, arquivos sensíveis expostos e vulnerabilidades conhecidas.
            Essa ferramenta auxilia na avaliação preliminar da segurança de aplicações web, permitindo detectar falhas 
            comuns que podem comprometer o ambiente\citep{moreno2019introduccao}.
        
        \subsection{Gobuster e Dirsearch}

            Ferramentas como Gobuster e Dirsearch são utilizadas para a enumeração de diretórios, arquivos e recursos ocultos em aplicações web.
            Essas ferramentas exploram listas de palavras para identificar caminhos não diretamente acessíveis, auxiliando na descoberta de superfícies
            de ataque adicionais. A enumeração detalhada de recursos expostos é uma etapa relevante na avaliação da 
            segurança de aplicações\citep{stallings2017network}.

        \subsection{SQLMap e Commix}
        
            O SQLMap é uma ferramenta destinada à identificação e exploração automatizada de vulnerabilidades de injeção SQL, enquanto o Commix
            foca na detecção de falhas de injeção de comandos no sistema operacional. Ambas reduzem o esforço manual do testador ao automatizar
            testes complexos, permitindo validar o impacto real das vulnerabilidades encontradas \citep{moreno2019introduccao}.

        \subsection{Metasploit}

            O Metasploit é uma plataforma abrangente para exploração de vulnerabilidades e realização de testes de pós-exploração.
            Ele permite avaliar as consequências práticas das falhas identificadas, auxiliando na compreensão do impacto sobre a
            confidencialidade, integridade e disponibilidade dos sistemas. Frameworks de exploração são 
            fundamentais para validar a efetividade dos mecanismos de segurança existentes\citep{stallings2017network}.

        \subsection{Netcat}

            O Netcat é uma ferramenta de rede de propósito geral utilizada para estabelecer conexões TCP/UDP, realizar testes de 
            conectividade e interagir com serviços de forma direta. Em atividades de pentest, é frequentemente empregado para validação
            de portas e serviços, simulação de comunicação cliente-servidor e apoio a testes que exigem troca de dados em baixo nível.
            Por permitir inspeção e interação com serviços de forma simples, o Netcat é útil em etapas de reconhecimento e validação de achados,
            complementando scanners automatizados e auxiliando na compreensão prática do comportamento de rede \citep{stallings2017network,moreno2019introduccao}.
        
        \subsection{Caido}

            O Caido é uma ferramenta utilizada como proxy de interceptação para análise de tráfego HTTP/HTTPS, permitindo a inspeção,
            modificação e repetição de requisições entre cliente e servidor. Esse tipo de ferramenta é amplamente empregado em testes
            de segurança de aplicações web, pois possibilita observar o funcionamento interno das comunicações e identificar falhas
            relacionadas a autenticação, autorização e validação de entradas. A análise do tráfego de aplicações web
            é fundamental para compreender o comportamento das requisições e detectar vulnerabilidades que não são facilmente identificadas
            por scanners automatizados \citep{moreno2019introduccao}.
            
            Ao atuar como intermediário entre o navegador e a aplicação alvo, o Caido permite que o testador avalie de forma prática
            o impacto de alterações em parâmetros e cabeçalhos HTTP, auxiliando na identificação de falhas lógicas e de configuração.
            Dessa forma, a utilização de ferramentas de interceptação complementa o uso de scanners e frameworks de exploração, 
            contribuindo para uma análise mais abrangente da segurança de aplicações web \citep{stallings2017network}.

\section{Virtualização e Orquestração (Docker)}

\section{Integração e Normalização de Saídas}

\section{Desenvolvimento da Interface Gráfica e Fluxos Guiados}

\section{Implementação de LLM para Suporte ao Usuário}

\section{Geração de Relatórios Técnicos}